% BHU PYQ -- Professional Exam Template
\documentclass[11pt,a4paper]{article}
\usepackage[a4paper,top=2.5cm,bottom=2.5cm,left=2.8cm,right=2.8cm,headheight=14pt]{geometry}
\usepackage{amsmath,amssymb}
\usepackage[T1]{fontenc}
\usepackage[utf8]{inputenc}
\usepackage{lmodern}
\usepackage{microtype}
\usepackage{fancyhdr}
\usepackage{enumitem}
\usepackage{hyperref}
\usepackage{lastpage}
\usepackage{parskip}

\hypersetup{colorlinks=true,urlcolor=black,linkcolor=black,
  pdftitle={CS-102: Digital Logics and Circuits -- BHU PYQ 2023-24}}

\pagestyle{fancy}\fancyhf{}
\renewcommand{\headrulewidth}{0.4pt}\renewcommand{\footrulewidth}{0.4pt}
\fancyhead[L]{\small   \textit{Banaras Hindu University}}
\fancyhead[R]{\small   \textit{CS-102: Digital Logics and Circuits}}
\fancyfoot[C]{\small Page  \thepage\ of \pageref{LastPage}}


\newlist{parts}{enumerate}{2}
\setlist[parts,1]{label=(\alph*),leftmargin=*,itemsep=4pt,topsep=3pt}
\setlist[parts,2]{label=(\roman*),leftmargin=*,itemsep=2pt,topsep=2pt}

\newcommand{\pts}[1]{\hfill   \textit{[#1]}}


\begin{document}

\begin{center}
  {\large\bfseries BANARAS HINDU UNIVERSITY}\\[2pt]
  {\normalsize B.Sc.\ (Hons.) Semester II Examination, 2023-24}\\[6pt]
  \rule{0.8\linewidth}{0.5pt}\\[6pt]
  {\large\bfseries Paper: CS-102 --- Digital Logics and Circuits}\\[4pt]
  {\normalsize Time: Three Hours \hfill Maximum Marks: 70}
\end{center}
\rule{\linewidth}{0.4pt}
\smallskip



\noindent   \textit{Instructions: Attempt   \textbf{five} questions including Question No.~1, which is compulsory. Write your Roll No.\ at the top immediately on receipt of this question paper.}
\bigskip




\noindent   \textbf{Question 1.}

\begin{parts}
  \item How does changing the base of a number system impact the representation and arithmetic operations performed on numbers?\pts{3}
  \item What is an unweighted and weighted code? Provide an example of each.\pts{3}
  \item What is the primary difference between combinational and sequential circuits?\pts{3}
  \item Explain the concept of a multiplexer and its applications. Design a 4-to-1 multiplexer using basic logic gates.\pts{3}
  \item How many 1's are present in the binary representation of $4 imes256 + 8 imes16 + 12$?\pts{2}
\end{parts}

\bigskip



\noindent   \textbf{Question 2.}

\begin{parts}
  \item Compute: $(B1.21)_{16} + (211)_{4} + (131)_{8} = (?)_{8}$\pts{4}
  \item How is a negative number represented in sign-magnitude form? Calculate the sum of $-36$ and $+52$ using signed number addition.\pts{5}
  \item Express the following functions in Sum of Minterms and Product of Maxterms and implement using a universal gate of your choice:\pts{5}
    
\begin{parts}
      \item $F(A,B,C,D) = A'D + C'D' + ABD$
      \item $F(A,B,C) = (ABC + BC)(A'B'C + BC)$
    \end{parts}
\end{parts}

\bigskip



\noindent   \textbf{Question 3.}

\begin{parts}
  \item Minimize the expression with don't-care conditions $F(A,B,C,D) = \sum(0,1,5,13,14,15) + \sum d(3,4,7,10,11)$ using K-map. Draw the logic diagram of the minimized circuit. Why are reflected codes used in K-Map?\pts{7}
  \item Determine the Prime-Implicants and essential Prime-Implicants of the following Boolean function using the Tabulation method:\pts{7}
    \[F(A,B,C,D,E,F,G) = \Sigma(20,28,30,52,60)\]
\end{parts}

\bigskip



\noindent   \textbf{Question 4.}

\begin{parts}
  \item What are Half and Full-Adder circuits? Explain in detail with a logic diagram and truth table. For the addition of two 8-bit numbers, what is the required number of half adders and full adders? Explain with a logic diagram.\pts{5}
  \item Design a circuit that compares two 4-bit numbers $X$ and $Y$. The outcome of the comparison is specified by three binary variables indicating whether $X=Y$, $X>Y$, or $X<Y$.\pts{5}
  \item Implement a Full-Subtractor circuit using a multiplexer.\pts{4}
\end{parts}

\bigskip



\noindent   \textbf{Question 5.}

\begin{parts}
  \item What is ROM? Design a combinational circuit using a ROM that accepts a 2-bit number and generates an output binary number equal to the cube of the input number.\pts{5}
  \item What is a code converter, and why is it important in digital systems? Design a BCD to Excess-3 code converter.\pts{5}
  \item What is a Programmable Logic Array (PLA)? Explain its basic structure and functionality.\pts{4}
\end{parts}

\bigskip



\noindent   \textbf{Question 6.}

\begin{parts}
  \item Explain the various types of flip-flops used in digital circuits: SR, JK, D, and T flip-flops. Explain how each can be implemented using basic logic gates. For each type, describe its truth table and characteristic equation.\pts{6}
  \item What is an excitation table for SR, D, T, and JK flip-flops, and how does it differ from a truth table?\pts{4}
  \item What is a Register? Design a 4-bit Register.\pts{4}
\end{parts}

\bigskip



\noindent   \textbf{Question 7.}

\begin{parts}
  \item What are counters and their types? Design a 3-bit synchronous up-counter using a T flip-flop. Explain with a timing diagram.\pts{7}
  \item What is the primary purpose of state reduction and state assignment in sequential circuit design? Explain with an example.\pts{7}
\end{parts}

\vfill\rule{\linewidth}{0.4pt}

\begin{center}\small
  \href{https://bhu-exam.vercel.app/}{Click here to visit our website}\\[4pt]
  \href{https://me-aryan.vercel.app/}{Click here to visit our contributor}\\[4pt]
  \href{https://quashan-me.vercel.app/}{Click here to visit our contributor}
\end{center}
\end{document}